\documentclass[11pt]{article}

\usepackage[utf8]{inputenc}
\usepackage[T1]{fontenc}
\usepackage[spanish]{babel}
\usepackage{amsmath,amssymb}
\usepackage{geometry}

\geometry{margin=2.8cm}

\begin{document}

\begin{center}
    {\Large \textbf{Plan de trabajo para monografía de grado}}\\[1em]
    {\large Micaela Berezan}\\[0.5em]
    Orienta:  Juliana Xavier 
\end{center}

\vspace{1em}

La presente tiene como objetivo presentar el plan de trabajo correspondiente a
la monografía de grado de la estudiante Micaela Berezan, así como establecer un
cronograma tentativo para su desarrollo y posterior defensa.

\vspace{1em}

\section*{Temática a abordar}

El objetivo de la monografía es estudiar diversas herramientas de topología
algebraica y topología geométrica vinculadas a la teoría de encajes de esferas,
con especial énfasis en la interacción entre resultados clásicos de homología y
ejemplos de encajes salvajes en  dimensión baja.\\



Como punto de partida se introducirán las nociones básicas de homología simplicial y singular, presentando las definiciones fundamentales, ejemplos y propiedades elementales necesarias para la construcción de los grupos de homología, así como el resultado que identifica el abelianizado del grupo fundamental con el primer grupo de homología para espacios conexos por caminos. A continuación se desarrollará la sucesión exacta de Mayer-Vietoris, obtenida a partir del lema de la serpiente aplicado a una sucesión exacta de complejos de cadenas. Esta herramienta permitirá calcular la homología de las esferas y servirá como ingrediente central en la demostración del teorema de separación de Jordan-Brouwer, que se estudiará a continuación.\\


Una vez desarrolladas estas herramientas se abordará el teorema generalizado de
Schoenflies, siguiendo la demostración de Morton Brown, que establece que, dado
un encaje de $S^{n}\times I$ en $S^{n+1}$, la clausura de cada componente conexa
del complemento es una $(n+1)$-bola.\\

Para completar el panorama de cómo se ven los encajes topológicos se presentarán resultados sobre la estructura del conjunto de puntos no localmente planos de encajes de codimensión uno en dimensión mayor que tres, en particular el teorema de Robion Kirby que describe las restricciones topológicas que debe satisfacer dicho conjunto cuando es no trivial.\\

Finalmente, se estudiarán diversos ejemplos clásicos de encajes de $S^{2}$ en
$S^{3}$, entre ellos la esfera de Alexander, la esfera de Antoine, la alfombra
de Bing y la esfera de Fox-Artin. Para cada uno de estos ejemplos se describirá
su construcción, se verificará que el espacio obtenido es una esfera topológica
y se analizarán las propiedades de sus regiones complementarias, así como la
naturaleza de sus puntos no localmente planos. Estos ejemplos ilustran los
fenómenos que distinguen a los encajes salvajes de los encajes localmente planos
y permiten apreciar el alcance de los resultados desarrollados previamente.\\

% Sugerimos el siguiente tribunal: Richard Mu\~niz, Aldo Portela y Rafael Potrie, Suplente: Juliana Xavier (orientadora).

\section*{Cronograma tentativo}

El siguiente cronograma es tentativo y podrá sufrir modificaciones durante el
desarrollo de la monografía.\\

Ya se comenz\'o a trabajar el a\~no pasado habiendo completado el estudio preliminar de la homología simplicial y singular, la sucesión de Mayer--Vietoris, el teorema de
    Jordan--Brouwer, la relación entre el grupo fundamental y el primer grupo de
    homología, y el teorema generalizado de Schoenflies.  Lo que sigue es el cronograma para este semestre que comienza y la posterior defensa a fin de a\~no:

\begin{itemize}
    \item \textbf{Agosto:}  Estudio de los resultados de Kirby sobre puntos no
    localmente planos y análisis de ejemplos clásicos de encajes salvajes:
    esfera de Alexander, esfera de Antoine, alfombra de Bing y esfera de
    Fox--Artin.


   \item \textbf{Setiembre-Octubre:} Redacción de la monografía.
   
    
\item \textbf{Noviembre:} Correcciones finales, entrega al tribunal y
    preparación de la defensa.
    
    \item \textbf{Diciembre:} Defensa.

    
\end{itemize}

\section*{Bibliografía principal}

\begin{itemize}
    \item A. Hatcher, \emph{Algebraic Topology}.
    % \item S. Mac Lane, \emph{Categories for the Working Mathematician}.
    \item M. Brown, \emph{A Proof of the Generalized Schoenflies Theorem}.
    \item R. Kirby, \emph{On the Set of Non-Locally Flat Points of a Submanifold of Codimension One}.
    \item R. Daverman, G. Venema, \emph{Embeddings in Manifolds}.
    \item R. H. Bing, \emph{A Wild Surface Each of Whose Arcs is Tame}.
\end{itemize}

\vspace{2cm}

\begin{flushright}
Micaela Berezan \hfill Orientadora 
\end{flushright}

\end{document}